\section{Lean 与形式化定理证明}

\subsection{Lean 简介}

Lean 是一种编程语言/定理证明器/交互式证明助手：
\begin{itemize}
    \item 定义自然数、加法等基本概念
    \item 声明并证明定理（如交换律 $a+b=b+a$）
    \item 证明以\textbf{证明树}表示：根节点为原始定理，每步分解为更简单的子目标
    \item 文件组织为项目，项目间可复用（类似软件开发中的库依赖）
\end{itemize}

\subsection{LeanDojo：开源定理证明基础设施}

\begin{importantbox}{LeanDojo（NeurIPS 2023）}
提供开源数据集、工具和模型：
\begin{itemize}
    \item 从人类编写的 Lean 代码中提取约 10,000 条定理和证明
    \item 提取每步的证明状态、使用的引理及其定义
    \item 训练 \textbf{ReProver}（检索增强证明器）：先检索相关引理，再生成下一步策略
\end{itemize}
\end{importantbox}

\subsection{Expert Iteration 改进}

用当前证明器尝试证明新定理 $\to$ 收集成功证明 $\to$ 加入训练集 $\to$ 训练更强证明器 $\to$ 迭代，直到性能饱和。

\subsection{定理证明的根本挑战}

\begin{warningbox}{无限动作空间}
围棋有 $19\times19$ 的有限棋盘，但数学证明的动作空间\textbf{实质无限}——难以仅靠人类数据覆盖，RL 探索也极其困难。
\end{warningbox}

\subsection{案例：不等式证明（LiPS）}

在数学奥林匹克不等式这一特定领域：
\begin{itemize}
    \item 将动作分为\textbf{缩放}（应用已知引理，有限可枚举）和\textbf{重写}（变换公式，无限空间交由 LLM 处理）
    \item 符号算法枚举所有缩放 + LLM 处理重写 + 启发式过滤
    \item O1 Preview 证明 0/20，O3 和 DeepSeek-R1 各证明 3--4/20，人类金牌选手 15/20，\textbf{LiPS 证明 16/20}
    \item 发现了人类未知的证明方法（仅用 AM-GM 配合巧妙重写）
\end{itemize}

